\documentclass{assignment}
\usepackage[utf8]{inputenc}
\usepackage[T1]{fontenc}
\usepackage{hyperref}

\lstset{
	inputencoding=utf8,
	literate=
	{á}{{\'a}}1
	{é}{{\'e}}1
	{í}{{\'i}}1
	{ó}{{\'o}}1
	{ú}{{\'u}}1
	{Á}{{\'A}}1
	{É}{{\'E}}1
	{Í}{{\'I}}1
	{Ó}{{\'O}}1
	{Ú}{{\'U}}1
	{ñ}{{\~n}}1
	{Ñ}{{\~N}}1
}
\begin{document}
	
	\assignmentTitle{Claudio Iván Esparza Castañeda}{Luis Alberto Balam Guzmán}{./Logo.png}{Maestría en Sistemas Embebidos}{Programación y Procesamiento de Datos}{1B. Práctica. Programación orientada a objetos en Python}
	
	\section*{Objetivo}
	Python, además de poder usarse con programación estructurada, es a su vez útil para la programación orientada a objetos (POO). \\
		La presente práctica busca establecer las bases de la programación orientada a objetos. se establece en primer lugar que la programación orientada a objetos permite trabajar con estructuras extensibles con propiedades que pueden usarse en copias de sí mismo.\\
		Este programa involucra el uso de un menú para el control del flujo, el cual debe ser robusto. Se cuenta con una clase padre "Sensor" que hereda sus atributos a las clases hijas, y en el proceso se sobreescriben los métodos de la clase padre.
	\section*{Pseudocódigo}
	\lstinputlisting[caption={Datos de sensores}, label={lst:sensores}]{Sensores.psc}
	\section*{Conclusiones}
	Se empleó una estructura basada en la tarea anterior para el menú, y mejorando la función de la validación de los datos flotantes. Se creó la clase padre, y además sus clases hijas, ajustando la lógica requerida. Soy consciente de que se puede entrar con mayor detalle a la POO y buentas prácticas a través de una mejor definición de los métodos setters y getters, o una definición general de los tipos de sensores y una sobreescritura usando una estructura for. Se decidió simplificar con el propósito de ajustarse a los requisitos propuestos para esta tarea.
	El repositorio de este proyecto es: \href{https://github.com/clivan/Programacion_y_Procesamiento_de_Datos_Infotec/1B}{Programación y procesamiento de Datos}.
	
\end{document}
